\section{Discussion}
\label{sec:discussion}

\subsection{Lift Is Necessary but Not Sufficient}

The Phase~9 caveat is methodological.  Lift over the copy baseline is a
useful diagnostic only when the representation remains discriminative.
If the online and target readouts both occupy tiny subspaces, a predictor
can achieve high diagonal cosine by mapping into the target subspace
without carrying much example-specific information.  The off-diagonal
cosine and KNN probe caught exactly this failure.  Future \CodeWM{}
claims should therefore report lift together with effective rank and KNN
retrieval.

\subsection{The Delta Is a Translation, Not a Vector Difference}

The failed state-space delta runs clarify what the model is learning.
The useful transition signal is not well represented as
$z_{t+1}-z_t$ in the final target space.  Instead, the predictor learns a
mapping from an online representation of the before-state, plus action
conditioning, into the target representation of the after-state.  When
\copycos{} becomes negative, this translation view becomes necessary:
copying the previous state is not a plausible explanation, but the
predictor can still align strongly with the target.  The caveat is that
translation must be discriminative.  VICReg is useful because it keeps
the readout high-rank while the predictor learns that translation.

\subsection{Deep Supervision and Shared Weights}

Deep supervision is especially natural for looped encoders.  In an
ordinary deep stack, auxiliary losses can be interpreted as helping early
layers receive stronger gradients, in the same way supervised contrastive
objectives do for image classifiers~\citep{khosla2020supcon}.  In a
weight-shared looped encoder, the interpretation is sharper: the same
block must learn a transformation that remains useful after multiple
recursion depths.  Predictive aux losses ask that shared block to
preserve usefulness at loop 1, loop 2, and the final loop, rather than
only at the endpoint.

This is a recipe rather than a one-off fix.  Shared weights make every
recursion depth reuse the same function, so a successful block cannot be
specialized only for a single layer index.  Auxiliary prediction losses
then prevent the easy failure mode: the network cannot defer all useful
geometry to the final readout while allowing intermediate states to
collapse.  However, deep supervision alone does not guarantee a healthy
pooled embedding.  The current recipe needs both pieces: aux prediction
for intermediate usefulness and VICReg for readout dimensionality.

\subsection{Connections to Temporal Straightening and Latent Reasoning}

Two concurrent lines of work offer complementary framings.
Concurrent work from the original JEPA
group~\citep{straightening2026} introduces a curvature regularizer for
JEPA-style latent planning that promotes temporally straight
trajectories: latent paths that can be linearly interpolated.  Our
VICReg objective acts as a discrete analogue: where they regularize
trajectory curvature continuously, we apply a per-loop variance hinge
that keeps the readout high-rank and isotropic, preventing the latent
trajectory from collapsing into a curved low-dimensional manifold.  However, the
per-loop straightness data (\S\ref{subsec:cross_seed}) show a
counterintuitive pattern: the peak checkpoint has \emph{lower}
straightness ($0.213$) than the post-fix run ($0.368$), yet performs
better on retrieval.  With only three loops yielding a single triple,
the straightness metric is not yet diagnostic; the connection remains
suggestive rather than confirmed.

\citet{jepareasoner2025} frame looped transformer blocks as latent
reasoning steps.  Under their framing, our predictive auxiliary
losses at intermediate loops are ``supervision on latent
chain-of-thought'' --- each loop must produce a representation useful
for the next reasoning step, not just the final answer.  This reframes
deep supervision from a gradient-flow technique to a latent-CoT
training signal, and suggests that the recipe may generalize to
non-code domains where looped encoders perform iterative refinement.

\subsection{Methodological Honesty: When Your Own Metric Lies}

The batch-128 KNN gallery we used during development has a random KNN@1
baseline of $3.9\%$, not the $0.02\%$ expected on the full validation
set.  Our screen-tier KNN@5 of $0.359$ corresponds to about $0.034$ on
full val; the absolute retrieval number was inflated by roughly
$12\times$.  Two features rescue the claim.  First, the VICReg-versus-baseline
\emph{ratio} transfers cleanly: $14\times$ KNN@5 and $25\times$ effective
rank on full val, slightly stronger than on the batch probe.  Second,
batch-128 remained directionally correct: every decision made against it
was corroborated by the full-val audit.  The recipe itself did not change
after the audit, only the reported numbers.  We report this inflation
explicitly rather than silently substituting the corrected numbers,
because the gameable-lift failure mode in Phase~9 Stage~2 makes the
discipline of reporting multiple metrics with their baselines load-bearing
for the paper.

\subsection{Where Scale Already Reshaped the Claim}

The $1000\times5000$ retrieval eval partially rewrote the external
claim.  On the development-scale $100\times300$ benchmark \CodeWM{}
matched or exceeded UnixCoder on every criterion, including the
strictest action-cosine threshold.  At paper scale the moderate and
R@1 wins survive but shrink, and on \code{cos@0.95} UnixCoder wins by
$0.009$ MRR.  We read this not as a failure of the mechanism but as a
healthy corrective on how the small-gallery benchmark was reading the
tail of the ranking distribution: edit retrieval at $1.3$M parameters
is \emph{competitive} with a $125.9$M-parameter general encoder,
not uniformly dominant.  The efficiency story ($100\times$ smaller,
$10\times$ lower CPU latency) is untouched.

\subsection{What Would Still Falsify the Story}

The mechanism claim (early-loop signal preservation plus VICReg readout
repair) survives the cross-seed audit in a narrower form than originally
hoped.  The step-28K peak \emph{did} partially disappear under other
seeds: cross-seed KNN@5 trails the peak by $2\text{--}3\times$
(\S\ref{subsec:cross_seed}).  What survives is the encoder-level claim:
VICReg reliably lifts encoder effective rank from $\sim\!2$ to
$53\text{--}64/128$ across all seeds.  The retrieval variance comes from
the predictor, whose effective rank ($5\text{--}6/128$) is uniformly
collapsed.  The story would be fully falsified if: (a)~the encoder
effective-rank improvement itself turns out to be seed-dependent rather
than consistent; (b)~the paper-grade \code{by\_joint} edge collapses
when CodeT5+ is added with its correct embedding head; or
(c)~predictor-specific regularization fails to close the
encoder--predictor rank gap, suggesting the bottleneck is architectural
rather than merely under-regularized.

\subsection{Predictor Collapse as the Next Bottleneck}

The cross-seed investigation surfaced a finding that was hidden by the
single-seed, lift-focused evaluation: the predictor is the bottleneck,
not the encoder.  Encoder effective rank is healthy ($53\text{--}64/128$)
across all seeds.  Predictor effective rank is $5\text{--}6/128$ across
\emph{all} checkpoints, including the peak.  VICReg was designed to fix
the encoder readout, and it succeeds at that.  But the predictor was
never explicitly regularized for rank, and the same collapse
pattern --- high cosine alignment, low cross-example discrimination ---
that previously affected the encoder readout now appears in the
predictor output.

This reframes the next axis of work.  Scaling data or encoder depth is
unlikely to help if the predictor compresses all encoder information into
a six-dimensional subspace.  Predictor-side VICReg, deeper predictor
architectures, or replacing the looped predictor with a single-pass
predictor are the natural next experiments.

\subsection{KernelWM}

\KernelWM{} is related but not identical.  Early results show high
prediction cosine and faster memorization, likely because the kernel
dataset is smaller than the CommitPackFT code-transition split.  For this
paper, \KernelWM{} should remain a secondary domain unless early-stopping
or data-augmentation checks show that the same geometry mechanism
generalizes without merely memorizing the smaller corpus.
